\chapter{Υποδείγματα ερωτήσεων και απαντήσεων}

Στο κεφάλαιο αυτό παρουσιάζονται screenshots από τα κύρια τμήματα της
εφαρμογής. Λειτουργεί ως σύντομο εγχειρίδιο χρήσης. Καθε σελίδα έχει σταθερό
\emph{PageHeader} με eyebrow, τίτλο και lede (italic), σταθερό
\emph{Footer} με credits και σταθερό \emph{Navbar} με τους 6 main routes
(Dashboard, Journals, Conferences, Authors, Years, Charts).

\section{Landing και dashboard}

\begin{figure}[H]
    \centering
    \includegraphics[width=0.95\linewidth,height=0.78\textheight,keepaspectratio]{screenshots/landing.png}
    \caption{Landing page: 3 KPI tiles (Articles, Authors, Venues) που γεμίζουν
    από το \texttt{GET /api/meta/totals}.}
\end{figure}

\begin{figure}[H]
    \centering
    \includegraphics[width=0.95\linewidth,height=0.78\textheight,keepaspectratio]{screenshots/dashboard.png}
    \caption{Dashboard με 4 KPI tiles, line chart ``Publications over time''
    και bar chart ``Top publishers by quartile''.}
\end{figure}

\section{List σελίδες}

\begin{figure}[H]
    \centering
    \includegraphics[width=0.95\linewidth,height=0.78\textheight,keepaspectratio]{screenshots/journals_list.png}
    \caption{Λίστα περιοδικών (\texttt{/journals}): server-side sort,
    sticky header, paginated table με filters για \texttt{best\_quartile},
    \texttt{best\_subject\_area} και \texttt{publisher}.}
\end{figure}

\begin{figure}[H]
    \centering
    \includegraphics[width=0.95\linewidth,height=0.78\textheight,keepaspectratio]{screenshots/conferences_list.png}
    \caption{Λίστα συνεδρίων (\texttt{/conferences}): filters για
    \texttt{rank\_value} και \texttt{primary\_for}, debounced text search,
    pagination και sortable columns.}
\end{figure}

\begin{figure}[H]
    \centering
    \includegraphics[width=0.95\linewidth,height=0.78\textheight,keepaspectratio]{screenshots/authors_list.png}
    \caption{Λίστα συγγραφέων (\texttt{/authors}): debounced
    \texttt{name\_query}, \texttt{min\_articles} filter, table με articles
    span. Default sort κατά \texttt{total\_articles} desc μέσω
    \texttt{materialized\_author\_profile}.}
\end{figure}

\begin{figure}[H]
    \centering
    \includegraphics[width=0.95\linewidth,height=0.78\textheight,keepaspectratio]{screenshots/years_list.png}
    \caption{Λίστα χρονιών (\texttt{/years}): line chart πάνω από τον
    πίνακα με 3 series (journal articles, conference articles, total)
    και sortable, linkable rows.}
\end{figure}

\section{Profile σελίδες}

\begin{figure}[H]
    \centering
    \includegraphics[width=0.95\linewidth,height=0.78\textheight,keepaspectratio]{screenshots/journal_profile.png}
    \caption{Profile περιοδικού (\texttt{/journals/\$id}): πλήρες ranking
    (best\_quartile, sjr\_index, h\_index), aggregate stats (total
    articles, distinct authors, avg articles/year), year range filter που
    \emph{επαναϋπολογίζει όλα τα stats}, line chart με 3 series ανά έτος,
    άρθρα σε paginated table.}
\end{figure}

\begin{figure}[H]
    \centering
    \includegraphics[width=0.95\linewidth,height=0.78\textheight,keepaspectratio]{screenshots/conference_profile.png}
    \caption{Profile συνεδρίου (\texttt{/conferences/\$id}): παρόμοιο layout
    με journal profile, εμπλουτισμένο με acronym, rank\_value, primary FoR.}
\end{figure}

\begin{figure}[H]
    \centering
    \includegraphics[width=0.95\linewidth,height=0.78\textheight,keepaspectratio]{screenshots/author_profile.png}
    \caption{Profile συγγραφέα (\texttt{/authors/\$id}): πρώτη / τελευταία
    χρονιά, σύνολο άρθρων, Μ.Ο. άρθρων ανά χρόνο, line chart, πλήρης
    βιβλιογραφία ως table με links σε venues και paper details.}
\end{figure}

\begin{figure}[H]
    \centering
    \includegraphics[width=0.95\linewidth,height=0.78\textheight,keepaspectratio]{screenshots/year_detail.png}
    \caption{Profile χρονιάς (\texttt{/years/\$year}): KPI tiles για
    τα 4 distinct counts, tabs για articles, journals, conferences,
    φίλτρα για conference, journal, author στο articles tab.}
\end{figure}

\section{Charts playground}

\begin{figure}[H]
    \centering
    \includegraphics[width=0.95\linewidth,height=0.78\textheight,keepaspectratio]{screenshots/charts_subject_area.png}
    \caption{Charts playground (\texttt{/charts}): line chart subject area
    yearly με checkbox multi-select στο sidebar, time-granularity selector,
    interactive legend, drag-to-zoom και reset.}
\end{figure}

\begin{figure}[H]
    \centering
    \includegraphics[width=0.95\linewidth,height=0.78\textheight,keepaspectratio]{screenshots/charts_publisher_quartile.png}
    \caption{Bar chart publisher quartile distribution: top-N publishers
    με 4 series Q1--Q4 stacked.}
\end{figure}

\begin{figure}[H]
    \centering
    \includegraphics[width=0.95\linewidth,height=0.78\textheight,keepaspectratio]{screenshots/charts_journal_metrics_scatter.png}
    \caption{Scatter plot journal metrics: X / Y / size pickers από 10 numeric
    metrics, color-by quartile / subject area / publisher, log/linear scale
    toggle per axis.}
\end{figure}

\begin{figure}[H]
    \centering
    \includegraphics[width=0.95\linewidth,height=0.78\textheight,keepaspectratio]{screenshots/charts_heatmap.png}
    \caption{Heatmap subject area $\times$ year: navy gradient που δείχνει
    article density ανά (area, year) bucket.}
\end{figure}

\begin{figure}[H]
    \centering
    \includegraphics[width=0.95\linewidth,height=0.78\textheight,keepaspectratio]{screenshots/charts_cumulative_growth.png}
    \caption{Cumulative growth: 3 running totals (journals, conferences,
    combined). Drag-to-zoom εμφανίζει badge με το zoomed εύρος.}
\end{figure}
